\documentclass[11pt]{article}
\usepackage[margin=1in]{geometry}
\usepackage{graphicx,url,longtable}
\title{An Atlas of NDM and OXA-48-like Carbapenemase Variation in Recent Outbreak Genomes: CARD-Validated Allele Calling, Structural Mapping of Variant Effects, and Porin-Loss Combination States}
\author{Builder 12 (ndm-oxa slice), amr-carbapenem-structure project\\Computational analysis of open public data; no wet-lab work}
\date{September 23, 2026}
\begin{document}
\maketitle
\begin{abstract}
Carbapenem-resistant Enterobacterales are driven largely by acquired carbapenemases, among which the metallo-beta-lactamase NDM and the class-D enzyme OXA-48 (with its OXA-48-like derivatives) dominate global surveillance. We built a validated, reproducible atlas of NDM and OXA-48-like variation across the NCBI Pathogen Detection resource (2.2 million isolate records), covering 91,774 blaNDM calls (60 alleles) and 41,865 OXA-48-like calls (31 alleles), with 28,289 isolates from 2023 onward joined to SNP-cluster outbreak identifiers. Our variant-calling pipeline (blastn against CARD) passed two locked positive-control gates: exact self-recovery of all 136 reference alleles (G1: 100\%), and 99.50\% agreement with NCBI AMRFinderPlus calls on 711 downloaded outbreak genomes (G2), with all eight residual disagreements resolved to reference-set divergence rather than caller error. Every variant position was mapped onto crystal structures of the matching enzyme class (NDM-1: PDB 4EYL with hydrolyzed meropenem; OXA-48: PDB 6P97 with imipenem), yielding 162 NDM and 68 OXA-48 variant positions annotated by structural context. Among 357 recent non-exact calls characterized at sequence level, 255 proved synonymous at protein level, while 51 carried amino-acid substitutions and 41 were internal-stop pseudogenes. The most striking novel observation is an NDM-5 variant carrying H250N, a substitution of a zinc-ligating residue (2.1 A in 4EYL), predicted to abolish metal-dependent catalysis. Porin analysis of 611 carbapenemase-positive K. pneumoniae and E. coli genomes found premature-stop disruption of OmpK35 in 170 of 309 Klebsiella genomes and OmpF stop mutations in 28 of 276 E. coli genomes; 5,803 isolates co-carry blaNDM and blaOXA-48-like genes. All gates were locked in writing before outcome data were examined; negative and thin results are preserved.
\end{abstract}

\section{Introduction}
Carbapenems are last-line beta-lactams for multidrug-resistant Gram-negative infection. Two carbapenemase families dominate the current epidemiology of Enterobacterales: the New Delhi metallo-beta-lactamase (NDM, class B) and the OXA-48-like oxacillinases (class D). NDM-1 was described in 2009 and has since diversified into dozens of named alleles, several with experimentally demonstrated effects on stability, zinc affinity, and substrate profile. OXA-48 and its derivatives (OXA-181, OXA-232, OXA-244 and others) hydrolyze carbapenems weakly but escape many phenotypic screens, and their spread is frequently amplified by porin loss (OmpK35/OmpK36 in Klebsiella pneumoniae, OmpC/OmpF in Escherichia coli), which raises effective carbapenem resistance by restricting drug entry.

Two gaps motivate this atlas. First, allele-resolved surveillance is fragmented: reference databases (CARD, ResFinder, NCBI's beta-lactamase allele registry) diverge in content and currency, and the allele composition of recent (2023 onward) outbreak genomes has not been summarized at the full scale of the NCBI Pathogen Detection resource with an explicitly validated caller. Second, variant interpretation is usually sequence-only: amino-acid changes are reported without verified mapping onto the corresponding enzyme structure, even though the structural consequences of a substitution at a zinc-ligating residue, an active-site loop, or the OXA-48 omega loop are qualitatively different from surface substitutions.

This study delivers: (i) a CARD-referenced, gate-validated allele caller with quantified agreement against NCBI AMRFinderPlus; (ii) the allele, species, geography, year, and SNP-cluster distribution of NDM and OXA-48-like carbapenemases across all of Pathogen Detection, with emphasis on 2023+ outbreak genomes; (iii) sequence-level characterization of all recent non-exact (potentially novel) calls, separating synonymous recovers from true amino-acid novelty; (iv) structural mapping of every variant position onto NDM-1 and OXA-48 crystal structures with ligand-context annotation; and (v) porin loss-of-function states in carbapenemase-positive genomes, quantifying combination genotypes.

\section{Methods}
\subsection{Success gates (locked before outcome data)}
Gates were locked in writing (repository: ndm-oxa/docs/GATES\_LOCKED.md, sha256 e74ab5a0...) before any outcome data were examined. G1: the caller must reproduce CARD reference alleles at 100\% exact recovery on CARD's own reference sequences. G2: on public genomes with independent NCBI AMRFinderPlus calls, agreement must be at least 95\% with disagreements resolved and counted. G3: every novel or outbreak-enriched variant claim must map onto the matching enzyme-class structure with alignment-verified residue correspondence; effect annotations are structural-context-based, never sequence-only. G4: porin loss-of-function calling must be validated on positive controls; combination claims are limited to co-occurrence quantification with verified/thin/missing counts. G5: verified/thin/missing counts in every table; negatives preserved.

\subsection{Reference data}
CARD (latest data tarball, 4.37 MB, sha256 94b89501...) supplied 6,059 protein-homolog-model nucleotide sequences. NDM alleles (n=69, names matching NDM-\textbackslash d+) and the OXA-48-like family were extracted. The OXA-48-like family was defined operationally, not from memory: all 1,274 CARD OXA proteins were aligned to OXA-48 (265 aa) with a global aligner, and members were taken at $\geq$90\% amino-acid identity, yielding 67 family alleles including OXA-48, OXA-163, OXA-181, OXA-232, OXA-244, OXA-245 and recent designations through OXA-1309. Borderline members (OXA-54, OXA-436, OXA-535, OXA-731 at 91--92\% identity) are disclosed. NCBI's beta-lactamase allele registry (Allele.tab, 1,557 entries) was downloaded as an independent cross-reference.

\subsection{Population data: NCBI Pathogen Detection}
The Pathogen Detection isolate backend (pathogens-srv, the solr API behind the public Isolates Browser) was queried directly for AMRFinderPlus records: all element\_symbol:blaNDM* calls (91,774) and all calls matching the 67 OXA-48-like symbols (72,691 raw, 41,865 after exact family-symbol filtering; the prefix query initially over-matched the unrelated Campylobacter-lineage alleles blaOXA-486/488/489, which were excluded). Collection date, geography, host, epi type, species, assembly and BioSample accessions, call method, percent identity/coverage, and reference-database version were retained. All 28,289 unique assemblies with collection dates in 2023 or later were joined to Pathogen Detection SNP-cluster (PDS) identifiers via the isolates collection (100-accession batches).

\subsection{Genome download and allele calling}
722 assemblies were selected for download: (a) all 324 recent assemblies with non-exact blaNDM calls (BLASTP/BLASTX/INTERNAL\_STOP/PARTIAL methods; contig-end truncations excluded as assembly artifacts), and (b) a reproducible pseudo-random (md5-keyed) sample of up to 25 recent K. pneumoniae and E. coli genomes per major NDM/OXA-48-like allele, for porin analysis. 711 of 722 were delivered by the NCBI Datasets v2 API (GENOME\_FASTA); 11 failed delivery and are counted as missing. Genomes were concatenated into 8 chunks of ~100 with accession-prefixed headers; a single blastn 2.17.0+ pass queried all 136 CARD references plus four porin references per chunk. A critical configuration bug was caught during validation: blastn -max\_target\_seqs truncates per query across the whole multi-genome database, silently dropping hits (47,756 hits at the default limit versus 113,357 at 20,000); all results use the corrected setting.

Allele calls: a gene copy is a cluster of hits on one contig; a copy is called exact when identity is 100\% and the alignment covers the full reference length, otherwise NONEXACT. G1 self-recovery: all 136 reference sequences recovered their own allele exactly (136/136, no ties). G2: against NCBI AMRFinderPlus calls on the 711 genomes, agreement was 1,599/1,607 (99.50\%). All eight disagreements were resolved: six involve alleles present in NCBI's reference gene catalog but absent from CARD (blaNDM-16b, blaNDM-94) or CARD/NCBI reference-sequence divergence (blaNDM-4); two are multi-copy genomes where NCBI reported one exact copy among several non-exact ones.

\subsection{Novel-variant characterization and structural mapping}
For every recent NONEXACT call, the hit region was extended to the full gene length on its contig, translated, and aligned to the closest CARD reference protein; substitutions, insertions/deletions, and internal stop codons were enumerated. Structural mapping: CARD NDM-1 and OXA-48 protein sequences were aligned to PDB chain sequences (NDM-1: 4EYL chain A, NDM-1 bound to hydrolyzed meropenem, 1.9 A, zinc present; OXA-48: 6P97 chain A, imipenem complex, 1.8 A; apo controls 3SPU and 4S2P). For each variant position, the minimum atomic distance to zinc ions (NDM), to the bound beta-lactam ligand, and to catalytic residues (OXA-48: S70, carbamylated K73) was computed from coordinates. Zones: zinc\_first\_shell ($\leq$4 A to Zn), active\_site\_near ($\leq$10 A to Zn), catalytic\_near ($\leq$8 A to S70/K73), ligand\_near ($\leq$8 A to ligand), surface\_other, or not-in-structure (signal peptide and uncovered termini).

\subsection{Reproducibility}
Every figure, table and count in this paper regenerates from the committed pipeline (ndm-oxa/code/) against the checksummed payloads listed in MANIFEST.sha256 (44 entries covering reference databases, raw API pulls, accession lists, PDB files, assembly archives and all intermediate tables). Gates were committed to version control before any results-bearing commit, so repository history independently proves gate precedence.
\subsection{Porin analysis}
References: OmpK35 (1,080 bp) and OmpK36 (1,119 bp) from K. pneumoniae AF21 (NZ\_OZ547068.1), OmpC (1,104 bp) and OmpF (1,089 bp) from E. coli K-12 MG1655 (NC\_000913.3). For each K. pneumoniae and E. coli genome, hits were extended to full gene length, translated, and classified: INTACT (full coverage, no internal stop, indel $<$6 bp), LOF\_STOP ($\geq$1 internal stop codon), INSERTION/DELETION ($\geq$6 bp indel), PARTIAL (coverage 60--90\%), MISSING. Positive controls (G4): synthetic spike-in mutants through the same code path (premature stop, loop-3 12-bp insertion, large deletion) were all detected with the expected class; intact references classified intact. A published porin-loss genome control remains pending and is disclosed as a gap. Caveat: indel calls against a single-strain reference are confounded by loop-length polymorphism and are reported as thin evidence; stop-codon calls are unambiguous.

\section{Results}
\subsection{Population atlas}
NDM: 91,774 calls, 60 distinct alleles. blaNDM-5 (42,537) and blaNDM-1 (39,463) dominate all-time; among 2023+ isolates blaNDM-5 leads blaNDM-1 (21,110 vs 13,496), consistent with continued displacement. K. pneumoniae (18,066) and E. coli (11,688) carry most recent NDM, followed by Enterobacter hormaechei (3,319) and Acinetobacter baumannii (2,167). blaNDM-5 spans 1,328 SNP clusters (largest: PDS000267523, 1,679 isolates); blaNDM-1 spans 1,134.
OXA-48-like: 41,865 calls, 31 alleles. blaOXA-48 (20,139 all-time; 5,208 recent) leads, with blaOXA-181 (10,707/3,791) and blaOXA-232 (6,713/1,539) established; blaOXA-244 shows fast recent growth (1,505 recent of 2,700 all-time). blaOXA-232 is unusually cluster-concentrated (top two clusters PDS000261518/19 hold 484 of 1,539 recent isolates), suggesting ongoing clonal outbreaks. Co-carriage of blaNDM plus blaOXA-48-like occurs in 5,803 assemblies.
\begin{table}[h]\centering\small\caption{Top blaNDM alleles in NCBI Pathogen Detection (per-isolate calls, deduplicated).}\label{tab:ndm}\begin{tabular}{lrr}\hline
\textbf{Allele}&\textbf{All-time}&\textbf{2023+}\\\hline
blaNDM-5&42,537&21,110\\
blaNDM-1&39,463&13,496\\
blaNDM-7&5,474&3,648\\
blaNDM-4&1,520&493\\
blaNDM&1,513&514\\
blaNDM-9&385&62\\
blaNDM-14&167&113\\
blaNDM-6&133&53\\
blaNDM-13&96&48\\
blaNDM-50&78&74\\
\hline\end{tabular}\end{table}
\begin{table}[h]\centering\small\caption{Top blaOXA-48-like alleles (family = >=90 percent amino-acid identity to OXA-48).}\label{tab:oxa}\begin{tabular}{lrr}\hline
\textbf{Allele}&\textbf{All-time}&\textbf{2023+}\\\hline
blaOXA-48&20,139&5,208\\
blaOXA-181&10,707&3,791\\
blaOXA-232&6,713&1,539\\
blaOXA-244&2,700&1,505\\
blaOXA-484&452&290\\
blaOXA-204&232&2\\
blaOXA-1207&229&200\\
blaOXA-1205&217&182\\
blaOXA-245&119&0\\
blaOXA-252&98&62\\
\hline\end{tabular}\end{table}
\begin{table}[h]\centering\small\caption{Porin status in carbapenemase-positive genomes (711-genome validation set). LOF stop = premature stop codon (verified-grade). Insertions/deletions (>=6 bp) are confounded by loop-length polymorphism vs the single-strain reference (thin evidence).}\label{tab:porin}\begin{tabular}{lrrrrrr}\hline
\textbf{Porin}&\textbf{Intact}&\textbf{Stop}&\textbf{Ins}&\textbf{Del}&\textbf{Miss/P}&\textbf{n}\\\hline
OmpK35 (Kp)&135&170&2&0&2&309\\
OmpK36 (Kp)&54&4&20&231&0&309\\
OmpF (Ec)&209&28&0&39&6&282\\
OmpC (Ec)&132&2&61&87&0&282\\
\hline\end{tabular}\end{table}

\begin{figure}[h]\centering\includegraphics[width=\textwidth]{fig1_allele_trends.png}\caption{Allele calls by collection year (NCBI Pathogen Detection).}\end{figure}

\subsection{Validation gates}
G1: 136/136 exact reference self-recovery. G2: 99.50\% (1,599/1,607) agreement with NCBI AMRFinderPlus; disagreements resolved to reference-set divergence. G4: spike-in positive controls all detected; published-genome control pending (disclosed). Gate evidence and the max\_target\_seqs defect are documented in the repository.

\subsection{Novel variants are mostly not novel at protein level}
Of 357 recent non-exact blaNDM calls characterized: 255 (71\%) are synonymous-only at protein level (NCBI bare ``blaNDM'' BLASTP calls at 99.6\% nucleotide identity encode exactly NDM-1/NDM-5 protein), 51 carry amino-acid substitutions, and 41 carry internal stop codons (pseudogenes). Nucleotide-level novelty therefore overstates protein-level novelty roughly three-fold in this call class.

\subsection{Structural mapping}
162 NDM variant positions: 1 zinc-first-shell, 22 active-site-near, 6 catalytic-near, 93 surface, 40 outside structure coverage (signal peptide/termini). 68 OXA-48-like positions: 6 catalytic-near, 11 ligand-near, 39 surface, 12 uncovered. Notable: omega-loop-adjacent R214S substitutions (OXA-1167, OXA-1201, OXA-1205) sit at the residue implicated in piperacillin substrate switching; S212 changes (OXA-163 S212D and others) neighbor the catalytic S70/K73 dyad. Two novel findings stand out: (i) GCA\_046367555.1 carries NDM-5 with H250N, replacing a Zn2-ligating histidine (2.1 A from the ion) - predicted catalytically dead, i.e., a circulating NDM pseudogene on an otherwise intact background; (ii) GCA\_059932675.2 carries NDM-1 with four clustered substitutions at the L10 loop edge (D212W, K214A, A215R, K216P), an unusual burden flagged for re-sequencing verification.
\begin{figure}[h]\centering\includegraphics[width=0.85\textwidth]{fig2_structure_zones.png}\caption{Structural context of variant positions.}\end{figure}

\subsection{Porin-loss combination states}
OmpK35 premature stops are common in carbapenemase-positive K. pneumoniae (170/309 genes; verified-grade). OmpF stops occur in 28/276 E. coli genes. OmpK36 indel calls (231 deletions, 20 insertions) are confounded by loop polymorphism and treated as thin; the 20 insertions are consistent with the known loop-3 insertion class. Across call sets, 56--82\% of genomes carry at least one porin disruption call (any class).
\begin{figure}[h]\centering\includegraphics[width=0.8\textwidth]{fig3_porin.png}\caption{Porin disruption rate by carbapenemase call set.}\end{figure}

\section{Discussion}
This atlas separates three quantities that surveillance summaries often conflate: nucleotide-level divergence (common, mostly synonymous), protein-level novelty (rare), and structurally consequential novelty (rarer still). The finding that 71\% of recent non-exact blaNDM calls encode known proteins shows how database-vintage artifacts inflate apparent diversity; conversely, the H250N zinc-ligand pseudogene shows that genuinely consequential variants do circulate and are recoverable only with structure-verified mapping. The continued rise of blaNDM-5 over blaNDM-1, the cluster-concentrated expansion of blaOXA-232, and the fast growth of blaOXA-244 are actionable surveillance signals. Porin disruption, especially OmpK35 stop mutations in Klebsiella, is the rule rather than the exception among carbapenemase-positive genomes, supporting combination-state reporting rather than enzyme-only calls.

Limitations: 711 of 722 target genomes were delivered; the OXA-48-like family boundary is a stated 90\% threshold; porin indel calls are thin due to loop polymorphism; the published-genome porin control is pending; MIC phenotypes are absent from Pathogen Detection metadata and are reported as missing, so no resistance-level claims are made.

\subsection{Per-gene porin biology}
The porin table reveals gene-specific disruption modes. OmpK35 loss in Klebsiella is dominated by premature stop codons (170 stop vs 135 intact), a verified-grade signal consistent with the dispensability of OmpK35 under carbapenemase pressure. OmpK36 shows few stops (4) but many indel-class calls (231 deletions, 20 insertions); because loop-length polymorphism between Klebsiella lineages inflates indel calls against any single reference, we do not interpret OmpK36 deletions as loss-of-function here, while the 20 insertions align with the documented loop-3 insertion class that constricts the channel. In E. coli, OmpF stops (28) exceed OmpC stops (2), mirroring the known OmpF-first loss order.
\subsection{Cluster epidemiology}
SNP-cluster joining converts allele counts into outbreak structure. blaNDM-5's largest cluster (PDS000267523, 1,679 isolates) is roughly twice the size of any blaNDM-1 cluster, and blaOXA-232's recent carriage is concentrated in two sibling clusters (PDS000261518/19), the signature of an expanding clone rather than dispersed plasmid spread. blaOXA-244's cluster count (139) relative to its recent total (1,505) indicates broad, multi-outbreak dissemination.
\section{Data and code availability}
Code, manifests (44 payloads with sha256), accession lists, tables and figures: repository amr-carbapenem-structure, ndm-oxa/ slice. All sources are open public data (CARD, NCBI Pathogen Detection, NCBI Datasets, RCSB PDB).

\begin{thebibliography}{9}\small
\bibitem{card2023} Alcock et al. CARD 2023. Nucleic Acids Res 2023. {\tt https://pmc.ncbi.nlm.nih.gov/articles/PMC9825576/}
\bibitem{ndmstability} Biochemical characterization of NDM variants reveals differences in protein stability. J Antimicrob Chemother. {\tt https://doi.org/10.1093/jac/dku403}
\bibitem{ndmevol} Evolution of NDM in the clinic: effects of NDM mutations on stability, zinc affinity, and mono-zinc activity. J Biol Chem. {\tt https://doi.org/10.1074/jbc.ra118.003835}
\bibitem{ndmstruct} King AM, Strynadka NCJ et al. Apo NDM-1 crystal structure (3SPU); NDM-1 with hydrolyzed meropenem (4EYL). RCSB PDB.
\bibitem{oxastruct} Structure, activity and thermostability investigations of OXA-163, OXA-181 and OXA-245. Acta Crystallogr D. {\tt https://doi.org/10.1107/s2053230x17013838}
\bibitem{oxareview} The Global Ascendency of OXA-48-Type Carbapenemases. Clin Microbiol Rev. {\tt https://doi.org/10.1128/cmr.00102-19}
\bibitem{r214} Improved hydrolysis of piperacillin by OXA-48-like R214G variants. Antimicrob Agents Chemother 2025. {\tt https://journals.asm.org/doi/10.1128/aac.01267-25}
\bibitem{porin2025} Proteogenomic analysis: blaOXA-48 copy numbers and OmpK36 loss in carbapenem resistance. Antimicrob Agents Chemother 2025. {\tt https://doi.org/10.1128/aac.00107-25}
\bibitem{loop3} Widespread emergence of OmpK36 loop 3 insertions. PLoS Pathog 2021. {\tt https://journals.plos.org/plospathogens/article?id=10.1371/journal.ppat.1010334}
\bibitem{pd} NCBI Pathogen Detection. {\tt https://www.ncbi.nlm.nih.gov/pathogens/}
\bibitem{oxa48struct} OXA-48 unbound (4S2P); OXA-48--imipenem complex (6P97). RCSB PDB.
\bibitem{oxa181} Characterization of OXA-181. Antimicrob Agents Chemother. {\tt https://pmc.ncbi.nlm.nih.gov/articles/PMC3186949/}
\end{thebibliography}

\appendix
\section{Complete NDM structural effect table}
All 162 variant positions observed across the 69 CARD NDM alleles and the recent novel calls, mapped to PDB 4EYL. dist-Zn = minimum atomic distance to a zinc ion; dist-lig = minimum distance to hydrolyzed meropenem.
\begin{longtable}{rrlrl}
\caption{NDM variant positions, structural annotation.}\label{tab:ndmfull}\\
\textbf{pos}&\textbf{pdb res}&\textbf{zone}&\textbf{d-Zn}&\textbf{examples}\\\hline
\endhead
1&-&NOT\_IN\_STRUCTURE(signal/uncovered)&-&GCA\_047893895.1:M1N; GCA\_047894095.1:M1N\\
2&-&NOT\_IN\_STRUCTURE(signal/uncovered)&-&GCA\_047893895.1:E2S; GCA\_047894095.1:E2S\\
3&-&NOT\_IN\_STRUCTURE(signal/uncovered)&-&GCA\_047893895.1:L3A; GCA\_047894095.1:L3A\\
4&-&NOT\_IN\_STRUCTURE(signal/uncovered)&-&GCA\_047919235.1:P4R; GCA\_049222765.1:P4R\\
5&-&NOT\_IN\_STRUCTURE(signal/uncovered)&-&GCA\_047919235.1:N5E; GCA\_049222765.1:N5E\\
6&-&NOT\_IN\_STRUCTURE(signal/uncovered)&-&GCA\_047919235.1:I6A; GCA\_049222765.1:I6A\\
7&-&NOT\_IN\_STRUCTURE(signal/uncovered)&-&GCA\_047919235.1:M7E; GCA\_049222765.1:M7E\\
8&-&NOT\_IN\_STRUCTURE(signal/uncovered)&-&GCA\_050379495.1:H8P; GCA\_054792955.1:H8A\\
9&-&NOT\_IN\_STRUCTURE(signal/uncovered)&-&GCA\_046368305.1:P9S; GCA\_050379495.1:P9V\\
10&-&NOT\_IN\_STRUCTURE(signal/uncovered)&-&GCA\_054792955.1:V10S; GCA\_054792955.1:V10S\\
11&-&NOT\_IN\_STRUCTURE(signal/uncovered)&-&GCA\_057961985.1:A11S\\
12&-&NOT\_IN\_STRUCTURE(signal/uncovered)&-&GCA\_038419775.1:K12R; GCA\_047919235.1:K12R\\
13&-&NOT\_IN\_STRUCTURE(signal/uncovered)&-&GCA\_038419775.1:L13I; GCA\_047919235.1:L13I\\
14&-&NOT\_IN\_STRUCTURE(signal/uncovered)&-&GCA\_054792955.1:S14D; GCA\_054792955.1:S14D\\
15&-&NOT\_IN\_STRUCTURE(signal/uncovered)&-&GCA\_038419775.1:T15R; GCA\_047919235.1:T15R\\
16&-&NOT\_IN\_STRUCTURE(signal/uncovered)&-&GCA\_038419775.1:A16C; GCA\_047919235.1:A16C\\
17&-&NOT\_IN\_STRUCTURE(signal/uncovered)&-&GCA\_038419775.1:L17I; GCA\_047919235.1:L17I\\
18&-&NOT\_IN\_STRUCTURE(signal/uncovered)&-&GCA\_038419775.1:A18D; GCA\_047919235.1:A18D\\
19&-&NOT\_IN\_STRUCTURE(signal/uncovered)&-&GCA\_054792955.1:A19T; GCA\_054792955.1:A19T\\
20&-&NOT\_IN\_STRUCTURE(signal/uncovered)&-&GCA\_038419775.1:A20E; GCA\_047919235.1:A20E\\
21&-&NOT\_IN\_STRUCTURE(signal/uncovered)&-&GCA\_038419775.1:L21R; GCA\_047919235.1:L21R\\
22&-&NOT\_IN\_STRUCTURE(signal/uncovered)&-&GCA\_038419775.1:M22V; GCA\_047919235.1:M22V\\
23&-&NOT\_IN\_STRUCTURE(signal/uncovered)&-&GCA\_038419775.1:L23H; GCA\_047919235.1:L23H\\
24&-&NOT\_IN\_STRUCTURE(signal/uncovered)&-&GCA\_038419775.1:S24A; GCA\_047919235.1:S24A\\
25&-&NOT\_IN\_STRUCTURE(signal/uncovered)&-&GCA\_038419775.1:G25R; GCA\_047919235.1:G25R\\
26&-&NOT\_IN\_STRUCTURE(signal/uncovered)&-&GCA\_057412325.1:C26N; GCA\_057412325.1:C26N\\
27&-&NOT\_IN\_STRUCTURE(signal/uncovered)&-&GCA\_038419775.1:M27N; GCA\_047919235.1:M27N\\
28&-&NOT\_IN\_STRUCTURE(signal/uncovered)&-&GCA\_057412325.1:P28L; GCA\_057412325.1:P28L\\
29&-&NOT\_IN\_STRUCTURE(signal/uncovered)&-&GCA\_037887805.1:G29D; GCA\_038419775.1:G29P\\
30&-&NOT\_IN\_STRUCTURE(signal/uncovered)&-&GCA\_038419775.1:E30D; GCA\_047919235.1:E30D\\
31&-&NOT\_IN\_STRUCTURE(signal/uncovered)&-&GCA\_038419775.1:I31D; GCA\_047919235.1:I31D\\
32&-&NOT\_IN\_STRUCTURE(signal/uncovered)&-&GCA\_038419775.1:R32W; GCA\_047919235.1:R32W\\
33&-&NOT\_IN\_STRUCTURE(signal/uncovered)&-&GCA\_057412325.1:P33S; GCA\_057412325.1:P33S\\
34&-&NOT\_IN\_STRUCTURE(signal/uncovered)&-&GCA\_038419775.1:T34A; GCA\_047919235.1:T34A\\
35&-&NOT\_IN\_STRUCTURE(signal/uncovered)&-&GCA\_038419775.1:I35N; GCA\_047919235.1:I35N\\
36&-&NOT\_IN\_STRUCTURE(signal/uncovered)&-&GCA\_053848185.1:G36S\\
37&-&NOT\_IN\_STRUCTURE(signal/uncovered)&-&GCA\_038419775.1:Q37N; GCA\_054792955.1:Q37N\\
38&-&NOT\_IN\_STRUCTURE(signal/uncovered)&-&GCA\_038419775.1:Q38W; GCA\_054792955.1:Q38W\\
39&-&NOT\_IN\_STRUCTURE(signal/uncovered)&-&GCA\_038419775.1:M39R\\
40&-&NOT\_IN\_STRUCTURE(signal/uncovered)&-&GCA\_038419775.1:E40P\\
43&43&surface\_other&17.4&GCA\_038419775.1:D43V\\
49&49&surface\_other&15.8&NDM-34:L49M\\
54&54&surface\_other&16.3&NDM-34:L54I\\
55&55&surface\_other&21.9&NDM-25:A55S\\
57&57&surface\_other&22.0&NDM-56:N57T\\
69&69&surface\_other&15.9&NDM-10:G69S; NDM-21:G69S\\
72&72&catalytic\_near&11.7&NDM-33:A72T; NDM-67:A72V\\
74&74&active\_site\_near&6.4&NDM-10:A74T\\
75&75&active\_site\_near&5.2&GCA\_054177645.1:S75R\\
76&76&active\_site\_near&7.3&GCA\_054177645.1:N76F\\
77&77&active\_site\_near&9.7&GCA\_054177645.1:G77Q\\
78&78&surface\_other&12.5&GCA\_054177645.1:L78R\\
79&79&surface\_other&16.0&NDM-66:I79L\\
80&80&surface\_other&18.6&GCA\_054177645.1:V80D\\
82&82&surface\_other&24.7&GCA\_054177645.1:D82Q\\
84&84&surface\_other&28.9&NDM-35:G84D\\
86&86&surface\_other&21.3&GCA\_054177645.1:V86P\\
87&87&surface\_other&18.4&NDM-56:L87R\\
88&88&surface\_other&15.1&NDM-17:V88L; NDM-20:V88L\\
89&89&surface\_other&11.7&GCA\_054177645.1:V89G\\
90&90&active\_site\_near&7.9&GCA\_054177645.1:D90G\\
91&91&active\_site\_near&9.6&GCA\_054177645.1:T91R\\
92&92&catalytic\_near&10.2&GCA\_054177645.1:A92Y\\
93&93&active\_site\_near&5.8&GCA\_054177645.1:W93R\\
94&94&surface\_other&12.2&GCA\_054177645.1:T94L\\
95&95&surface\_other&15.4&NDM-13:D95N; NDM-27:D95N\\
96&96&surface\_other&17.5&GCA\_054177645.1:D96R\\
98&98&surface\_other&13.9&GCA\_054177645.1:T98P\\
99&99&surface\_other&18.0&NDM-55:A99T\\
100&100&surface\_other&18.7&GCA\_054177645.1:Q100R\\
101&101&surface\_other&13.9&NDM-23:I101L\\
102&102&surface\_other&19.4&GCA\_054177645.1:L102D\\
103&103&surface\_other&22.0&GCA\_054177645.1:N103P\\
104&104&surface\_other&20.6&GCA\_054177645.1:W104Q; GCA\_059621975.1:W104R\\
105&105&surface\_other&19.6&GCA\_054177645.1:I105L\\
106&106&surface\_other&24.5&GCA\_054177645.1:K106D\\
108&108&surface\_other&25.0&GCA\_059020645.1:E108V\\
109&109&surface\_other&23.1&NDM-68:I109L\\
110&110&surface\_other&27.3&NDM-71:N110D\\
123&123&active\_site\_near&6.2&GCA\_056740755.1:Q123K; GCA\_059682985.1:Q123K\\
130&130&surface\_other&17.1&NDM-14:D130G; NDM-19:D130N\\
137&137&surface\_other&20.7&GCA\_054488385.1:I137V\\
147&147&surface\_other&14.4&GCA\_053513495.1:Q147G\\
148&148&surface\_other&10.0&GCA\_053513495.1:L148I\\
149&149&surface\_other&12.2&GCA\_053513495.1:A149G\\
150&150&surface\_other&14.2&GCA\_053513495.1:P150S; GCA\_059621975.1:P150R\\
151&151&surface\_other&12.5&NDM-47:Q151L\\
152&152&active\_site\_near&7.8&NDM-74:E152K; NDM-9:E152K\\
153&153&surface\_other&12.0&GCA\_053513495.1:G153H\\
154&154&active\_site\_near&8.8&NDM-11:M154V; NDM-12:M154L\\
155&155&surface\_other&13.7&GCA\_032868935.2:V155G; GCA\_059621975.1:V155G\\
156&156&surface\_other&15.2&GCA\_053513495.1:A156F\\
158&158&surface\_other&18.5&GCA\_053513495.1:Q158P; GCA\_061007145.1:Q158H\\
159&159&surface\_other&17.6&GCA\_059621975.1:H159D\\
160&160&surface\_other&18.2&GCA\_053513495.1:S160P\\
161&161&surface\_other&16.0&GCA\_053513495.1:L161K; GCA\_059621975.1:L161P\\
162&162&surface\_other&19.1&GCA\_051205915.1:T162I; GCA\_053513495.1:T162Y\\
163&163&surface\_other&15.6&GCA\_053513495.1:F163L\\
164&164&surface\_other&19.7&GCA\_053513495.1:A164T\\
167&167&surface\_other&18.2&GCA\_053513495.1:G167S\\
168&168&surface\_other&20.5&GCA\_053513495.1:W168L\\
169&169&surface\_other&19.5&GCA\_053513495.1:V169T\\
170&170&surface\_other&21.8&NDM-17:E170K\\
171&171&surface\_other&24.6&NDM-31:P171T\\
172&172&surface\_other&25.2&GCA\_053513495.1:A172P\\
173&173&surface\_other&21.2&GCA\_053513495.1:T173L\\
174&174&surface\_other&21.3&GCA\_053513495.1:A174K\\
176&176&surface\_other&24.9&GCA\_053513495.1:N176Q\\
178&178&surface\_other&24.6&GCA\_053513495.1:G178L\\
179&179&surface\_other&22.0&NDM-49:P179R\\
180&180&surface\_other&18.6&GCA\_053513495.1:L180V\\
181&181&surface\_other&18.2&GCA\_053513495.1:K181Q\\
185&185&surface\_other&12.6&NDM-58:P185S\\
200&200&surface\_other&25.3&NDM-10:G200R\\
202&202&surface\_other&20.4&NDM-60:D202N\\
206&206&active\_site\_near&8.7&NDM-70:G206A\\
212&212&catalytic\_near&11.2&GCA\_059932675.2:D212W\\
213&213&surface\_other&14.9&GCA\_059932675.2:S213T\\
214&214&catalytic\_near&15.5&GCA\_059932675.2:K214A\\
215&215&catalytic\_near&12.1&GCA\_059932675.2:A215R\\
216&216&catalytic\_near&14.0&GCA\_059932675.2:K216P\\
218&218&active\_site\_near&8.0&GCA\_059932675.2:L218R\\
219&219&active\_site\_near&8.0&GCA\_059932675.2:G219S\\
220&220&active\_site\_near&5.2&GCA\_046367555.1:N220X; GCA\_059932675.2:N220A\\
221&221&active\_site\_near&8.2&GCA\_046367555.1:L221X; GCA\_059932675.2:L221I\\
222&222&active\_site\_near&8.6&NDM-12:G222D; NDM-26:G222S\\
223&223&active\_site\_near&6.4&NDM-30:D223Y\\
224&224&active\_site\_near&8.4&GCA\_046367555.1:A224R; GCA\_059932675.2:A224M\\
225&225&surface\_other&10.6&GCA\_046367555.1:D225T; GCA\_050724915.1:D225E\\
226&226&surface\_other&13.5&GCA\_046367555.1:T226F\\
227&227&surface\_other&11.1&GCA\_046367555.1:E227P; GCA\_059932675.2:E227L\\
228&228&active\_site\_near&9.7&GCA\_046367555.1:H228C; GCA\_059932675.2:H228S\\
229&229&active\_site\_near&8.6&GCA\_046367555.1:Y229V; GCA\_059932675.2:Y229T\\
230&230&surface\_other&15.2&GCA\_059932675.2:A230T\\
231&231&surface\_other&16.9&GCA\_046367555.1:A231Q; GCA\_059932675.2:A231P\\
232&232&surface\_other&13.5&GCA\_046367555.1:S232I; GCA\_059932675.2:S232R\\
233&233&surface\_other&14.4&NDM-15:A233V; NDM-19:A233V\\
235&235&surface\_other&17.9&GCA\_046367555.1:A235E\\
236&236&surface\_other&14.1&GCA\_046367555.1:F236N\\
238&238&surface\_other&20.6&GCA\_046367555.1:A238D; GCA\_059932675.2:A238R\\
239&239&surface\_other&19.5&GCA\_046367555.1:A239N; GCA\_059932675.2:A239L\\
240&240&surface\_other&17.3&GCA\_046367555.1:F240V; GCA\_059932675.2:F240V\\
241&241&surface\_other&22.6&GCA\_046367555.1:P241I; GCA\_059932675.2:P241R\\
242&242&surface\_other&23.4&NDM-56:K242Q\\
243&243&surface\_other&19.8&GCA\_046367555.1:A243K; GCA\_052594815.1:A243T\\
244&244&surface\_other&20.2&GCA\_046367555.1:S244W; GCA\_059932675.2:S244P\\
245&245&surface\_other&17.1&GCA\_046367555.1:M245L; GCA\_055434245.1:M245I\\
246&246&surface\_other&13.7&GCA\_046367555.1:I246R\\
247&247&surface\_other&10.9&GCA\_046367555.1:V247L\\
248&248&active\_site\_near&7.8&NDM-22:M248L\\
249&249&active\_site\_near&4.2&GCA\_040442885.1:S249G; GCA\_046367555.1:S249Q\\
250&250&zinc\_first\_shell&2.1&GCA\_046367555.1:H250N\\
252&252&surface\_other&12.3&NDM-51:A252D\\
254&254&surface\_other&15.3&NDM-43:D254N; NDM-56:D254E\\
261&261&surface\_other&15.4&NDM-36:H261Y; NDM-37:H261Y\\
263&263&surface\_other&15.8&GCA\_059990755.1:A263V\\
264&264&surface\_other&18.0&GCA\_058143285.1:R264S\\
265&265&surface\_other&13.2&GCA\_050722475.1:M265I\\
266&266&surface\_other&14.6&NDM-28:A266V\\
267&267&surface\_other&17.7&NDM-43:D267N; NDM-44:D267N\\
268&268&surface\_other&19.4&NDM-56:K268T\\
270&270&surface\_other&15.7&NDM-20:R270H\\
\hline\end{longtable}
\section{Complete OXA-48-like structural effect table}
All 68 variant positions across the 67 CARD OXA-48-like alleles, mapped to PDB 6P97 (imipenem complex). d-S70/d-K73 = distance to catalytic residues; d-lig = distance to imipenem.
\begin{longtable}{rrlrl}
\caption{OXA-48-like variant positions, structural annotation.}\\
\textbf{pos}&\textbf{pdb res}&\textbf{zone}&\textbf{d-S70}&\textbf{examples}\\\hline
\endhead
3&-&NOT\_IN\_STRUCTURE&-&\\
5&-&NOT\_IN\_STRUCTURE&-&\\
8&-&NOT\_IN\_STRUCTURE&-&\\
10&-&NOT\_IN\_STRUCTURE&-&\\
11&-&NOT\_IN\_STRUCTURE&-&\\
13&-&NOT\_IN\_STRUCTURE&-&\\
14&-&NOT\_IN\_STRUCTURE&-&\\
15&-&NOT\_IN\_STRUCTURE&-&\\
16&-&NOT\_IN\_STRUCTURE&-&\\
17&-&NOT\_IN\_STRUCTURE&-&\\
21&-&NOT\_IN\_STRUCTURE&-&\\
22&-&NOT\_IN\_STRUCTURE&-&\\
23&23&surface\_other&24.2&OXA-54:K23N; OXA-547:K23R\\
26&26&surface\_other&25.2&OXA-731:Q26K\\
28&28&surface\_other&25.4&OXA-54:N28K; OXA-731:N28T\\
29&29&surface\_other&27.0&OXA-54:K29P; OXA-731:K29Q\\
33&33&ligand\_near&25.5&OXA-1307:A33S; OXA-54:A33T\\
34&34&ligand\_near&22.3&OXA-923:H34Y; OXA-934:H34Y\\
36&36&surface\_other&23.9&OXA-436:T36S; OXA-535:T36S\\
38&38&ligand\_near&21.4&OXA-199:H38Y\\
40&40&surface\_other&19.0&OXA-436:S40T; OXA-535:S40T\\
44&44&surface\_other&17.0&OXA-54:V44I\\
45&45&surface\_other&16.5&OXA-199:V45A; OXA-252:V45A\\
51&51&surface\_other&27.3&OXA-436:K51T; OXA-535:K51T\\
55&55&surface\_other&20.8&OXA-731:F55Y\\
58&58&surface\_other&20.7&OXA-436:N58D; OXA-535:N58D\\
59&59&surface\_other&21.5&OXA-1309:L59I\\
65&65&surface\_other&12.9&OXA-1308:A65V\\
75&75&catalytic\_near&6.8&OXA-918:P75T\\
98&98&surface\_other&15.9&OXA-204:Q98H\\
99&99&surface\_other&16.3&OXA-204:T99R\\
104&104&ligand\_near&11.4&OXA-1119:T104A; OXA-1181:T104A\\
110&110&surface\_other&12.8&OXA-1119:N110D; OXA-1181:N110D\\
120&120&catalytic\_near&4.2&OXA-519:V120L; OXA-788:V120L\\
121&121&catalytic\_near&8.6&OXA-1304:P121A\\
123&123&catalytic\_near&7.1&OXA-439:Y123H\\
125&125&surface\_other&12.4&OXA-245:E125Y; OXA-535:E125G\\
127&127&surface\_other&11.9&OXA-1242:A127S\\
129&129&surface\_other&16.9&OXA-731:Q129K\\
132&132&surface\_other&17.3&OXA-54:E132Q; OXA-731:E132Q\\
135&135&surface\_other&12.1&OXA-1305:M135I\\
136&136&surface\_other&16.9&OXA-731:S136N\\
141&141&surface\_other&18.0&OXA-566:A141D\\
153&153&surface\_other&12.6&OXA-416:V153L; OXA-436:V153L\\
154&154&surface\_other&10.7&OXA-199:D154G\\
155&155&surface\_other&9.3&OXA-1038:S155T; OXA-416:S155T\\
159&159&surface\_other&9.2&OXA-933:D159N\\
167&167&surface\_other&15.1&OXA-1038:T167I; OXA-1226:T167I\\
168&168&surface\_other&13.3&OXA-1119:E168Q; OXA-1181:E168Q\\
170&170&surface\_other&13.4&OXA-54:I170V\\
171&171&surface\_other&15.4&OXA-1039:S171T; OXA-1119:S171A\\
201&201&surface\_other&25.8&OXA-416:G201A; OXA-436:G201A\\
211&211&catalytic\_near&4.6&OXA-247:Y211S\\
212&212&catalytic\_near&6.1&OXA-1211:S212L; OXA-1240:S212G\\
213&213&ligand\_near&8.2&OXA-1119:T213I; OXA-162:T213A\\
214&214&ligand\_near&9.2&OXA-1167:R214S; OXA-1201:R214S\\
215&215&surface\_other&11.4&OXA-438:I215Y\\
216&216&surface\_other&13.8&OXA-438:E216D; OXA-517:E216K\\
217&217&surface\_other&11.6&OXA-1012:P217S; OXA-1055:P217S\\
218&218&surface\_other&11.6&OXA-1212:K218N; OXA-54:K218Q\\
219&219&ligand\_near&8.2&OXA-538:I219F\\
226&226&surface\_other&16.8&OXA-436:V226I\\
237&237&surface\_other&10.7&OXA-436:M237T; OXA-535:M237T\\
244&244&ligand\_near&12.5&OXA-1181:S244W; OXA-1205:S244W\\
245&245&ligand\_near&15.2&OXA-436:D245E; OXA-535:D245E\\
252&252&ligand\_near&16.3&OXA-436:A252T; OXA-535:A252S\\
256&256&ligand\_near&20.8&OXA-436:E256A; OXA-535:E256A\\
259&259&surface\_other&23.3&OXA-833:K259T\\
\hline\end{longtable}
\section{Novel variant calls with amino-acid changes or stops (2023+)}
92 of 357 characterized recent non-exact calls alter protein sequence or length; the remaining 265 are synonymous at protein level or partial.
\begin{longtable}{llll}
\caption{Protein-altering novel NDM/OXA-48-like calls.}\\
\textbf{assembly}&\textbf{closest ref}&\textbf{changes (n)}&\textbf{stops}\\\hline
\endhead
GCA\_031038655.3&NDM-3&E152K (1)&0\\
GCA\_031056435.1&NDM-1&D130N M154L (2)&0\\
GCA\_032671545.1&NDM-1&del257A del258A del259I del260T del261H del262T  (15)&0\\
GCA\_032868935.2&NDM-5&V155G (1)&0\\
GCA\_037887805.1&NDM-5&G29D (1)&0\\
GCA\_038419775.1&NDM-1&ins10S ins10E ins11E ins11H K12R L13I T15R A16C  (226)&6\\
GCA\_038487205.1&NDM-1&M154L (1)&0\\
GCA\_040442885.1&NDM-4&S249G (1)&0\\
GCA\_040442885.1&OXA-181&ins0L ins0N M1Y V3D L4E A5V A8C V9A ... (213)&12\\
GCA\_042177995.1&NDM-1&V88L M154L (2)&0\\
GCA\_044895485.1&NDM-1&V88L M154L (2)&0\\
GCA\_045445595.1&NDM-1&E152G (1)&0\\
GCA\_045445635.1&NDM-1&E152G (1)&0\\
GCA\_045445675.1&NDM-1&E152G (1)&0\\
GCA\_045445755.1&NDM-1&E152G (1)&0\\
GCA\_045718445.1&NDM-5&del256R del257A del258A del259I del260T del261H  (16)&0\\
GCA\_046367555.1&NDM-5&del219G N220X L221X G222X D223X A224R D225T T226 (48)&0\\
GCA\_046368305.1&NDM-5&P9S (1)&0\\
GCA\_047893855.1&OXA-232&A104T D110N Q168E A171S S214R (5)&0\\
GCA\_047893895.1&NDM-4&M1N E2S L3A (3)&0\\
GCA\_047894095.1&NDM-4&M1N E2S L3A (3)&0\\
GCA\_047894795.1&NDM-4&M1N E2S L3A (3)&0\\
GCA\_047919235.1&NDM-1&M1A E2P L3G P4R N5E I6A M7E del9P ... (232)&6\\
GCA\_048020515.1&NDM-4&M1N E2S L3A (3)&0\\
GCA\_048247675.1&NDM-1&V88L M154L (2)&0\\
GCA\_049222765.1&NDM-1&M1A E2P L3G P4R N5E I6A M7E del9P ... (232)&6\\
GCA\_049232105.1&NDM-1&M1A E2P L3G P4R N5E I6A M7E del9P ... (232)&6\\
GCA\_049236645.1&NDM-1&M1A E2P L3G P4R N5E I6A M7E del9P ... (232)&6\\
GCA\_049569975.1&NDM-1&V88L M154L (2)&0\\
GCA\_050379495.1&NDM-1&H8P P9V (2)&0\\
GCA\_050582225.1&NDM-5&del262T del263A del264R del265M del266A del267D  (10)&0\\
GCA\_050595725.1&NDM-1&V88L M154L (2)&0\\
GCA\_050595745.1&NDM-4&V88L (1)&0\\
GCA\_050690955.1&NDM-5&del262T del263A del264R del265M del266A del267D  (10)&0\\
GCA\_050722475.1&NDM-5&M265I (1)&0\\
GCA\_050724915.1&NDM-5&D225E (1)&0\\
GCA\_051205915.1&NDM-5&T162I (1)&0\\
GCA\_051414975.1&NDM-1&M1A E2P L3G P4R N5E I6A M7E del9P ... (232)&6\\
GCA\_052445325.1&NDM-4&V88L (1)&0\\
GCA\_052568655.1&NDM-1&E2D (1)&0\\
GCA\_052594815.1&NDM-5&A243T (1)&0\\
GCA\_052942835.1&NDM-1&M1A E2P L3G P4R N5E I6A M7E del9P ... (232)&6\\
GCA\_052942855.1&NDM-1&M1A E2P L3G P4R N5E I6A M7E del9P ... (232)&6\\
GCA\_052942875.1&NDM-1&M1A E2P L3G P4R N5E I6A M7E del9P ... (232)&6\\
GCA\_052942895.1&NDM-1&M1A E2P L3G P4R N5E I6A M7E del9P ... (232)&6\\
GCA\_052942915.1&NDM-1&M1A E2P L3G P4R N5E I6A M7E del9P ... (232)&6\\
GCA\_052942935.1&NDM-1&M1A E2P L3G P4R N5E I6A M7E del9P ... (232)&6\\
GCA\_052942955.1&NDM-1&M1A E2P L3G P4R N5E I6A M7E del9P ... (232)&6\\
GCA\_052942975.1&NDM-1&M1A E2P L3G P4R N5E I6A M7E del9P ... (232)&6\\
GCA\_052942995.1&NDM-1&M1A E2P L3G P4R N5E I6A M7E del9P ... (232)&6\\
GCA\_052943015.1&NDM-1&M1A E2P L3G P4R N5E I6A M7E del9P ... (232)&6\\
GCA\_052943055.1&NDM-1&M1A E2P L3G P4R N5E I6A M7E del9P ... (232)&6\\
GCA\_053194025.1&NDM-5&G29D (1)&0\\
GCA\_053336735.1&NDM-1&M1A E2P L3G P4R N5E I6A M7E del9P ... (232)&6\\
GCA\_053336755.1&NDM-1&M1A E2P L3G P4R N5E I6A M7E del9P ... (232)&6\\
GCA\_053337215.1&NDM-5&M1A E2P L3G P4R N5E I6A M7E del9P ... (231)&6\\
GCA\_053366035.1&NDM-1&E2D (1)&0\\
GCA\_053513495.1&NDM-5&del146N Q147G L148I A149G P150S Q151I E152N G153 (112)&8\\
GCA\_053848185.1&NDM-5&G36S (1)&0\\
GCA\_053848345.1&NDM-1&M1A E2P L3G P4R N5E I6A M7E del9P ... (232)&6\\
GCA\_054177645.1&NDM-1&del73V A74C S75R N76F G77Q L78R I79F V80D ... (175)&5\\
GCA\_054488385.1&NDM-1&I137V (1)&0\\
GCA\_054792955.1&NDM-1&del1M E2* L3M P4G N5S I6C M7A H8A ... (231)&6\\
GCA\_054792955.1&NDM-5&del1M E2* L3M P4G N5S I6C M7A H8A ... (229)&6\\
GCA\_054794435.1&NDM-71&D110N (1)&0\\
GCA\_055434245.1&NDM-5&M245I (1)&0\\
GCA\_056740755.1&NDM-1&Q123K M154L (2)&0\\
GCA\_056781405.1&NDM-1&M1A E2P L3G P4R N5E I6A M7E del9P ... (232)&6\\
GCA\_056969255.1&NDM-4&D130N (1)&0\\
GCA\_057174215.1&NDM-1&M1A E2P L3G P4R N5E I6A M7E del9P ... (232)&6\\
GCA\_057412325.1&NDM-1&M1* E2P ins3H ins3* ins3C P4* N5A I6G ... (228)&10\\
GCA\_057412325.1&NDM-5&M1* E2P ins3H ins3* ins3C P4* N5A I6G ... (228)&10\\
GCA\_057731325.1&NDM-5&del1M E2H L3* del5N I6L M7H H8* P9C ... (230)&10\\
GCA\_057961985.1&NDM-5&M1K E2C L3R N5L I6N M7G H8F P9M ... (28)&2\\
GCA\_058143285.1&NDM-5&R264S (1)&0\\
GCA\_058143815.1&NDM-1&V88L M154L (2)&0\\
GCA\_058195915.1&NDM-5&D130G (1)&0\\
GCA\_058313035.1&NDM-1&M1A E2P L3G P4R N5E I6A M7E del9P ... (232)&6\\
GCA\_058666915.1&NDM-1&M1A E2P L3G P4R N5E I6A M7E del9P ... (232)&6\\
GCA\_059020645.1&NDM-5&E108V (1)&0\\
GCA\_059034205.1&NDM-5&D130G (1)&0\\
GCA\_059621975.1&NDM-1&W104R A149* P150R Q151G E152D M154C V155G del157 (11)&1\\
GCA\_059637815.1&NDM-1&M1A E2P L3G P4R N5E I6A M7E del9P ... (232)&6\\
GCA\_059637835.1&NDM-1&M1A E2P L3G P4R N5E I6A M7E del9P ... (232)&6\\
GCA\_059682985.1&NDM-1&Q123K M154L (2)&0\\
GCA\_059932675.2&NDM-1&del211K D212W S213T K214A A215R K216P L218R G219 (52)&3\\
GCA\_059990755.1&NDM-1&A263V (1)&0\\
GCA\_061007145.1&NDM-1&Q158H (1)&0\\
GCA\_061012805.1&NDM-1&M1A E2P L3G P4R N5E I6A M7E del9P ... (232)&6\\
GCA\_061073645.1&NDM-1&M1A E2P L3G P4R N5E I6A M7E del9P ... (232)&6\\
GCA\_061369735.1&NDM-5&del245M del246I del247V del248M del249S del250H  (27)&0\\
GCA\_977071415.1&NDM-1&M1A E2P L3G P4R N5E I6A M7E del9P ... (232)&6\\
\hline\end{longtable}

\section{Geographic and host-species distribution (2023+)}
\begin{table}[h]\centering\small\caption{Top countries of isolation, recent isolates.}\begin{tabular}{lrlr}\hline
\textbf{NDM country}&\textbf{n}&\textbf{OXA-48-like country}&\textbf{n}\\\hline
USA&26,557&USA&4,203\\
China&3,073&Slovenia&1,237\\
Netherlands&1,252&Netherlands&1,119\\
India&1,091&South Africa&692\\
Slovakia&758&Denmark&632\\
Canada&721&Spain&600\\
Italy&661&India&467\\
Slovenia&594&Slovakia&450\\
Germany&437&Thailand&323\\
France&328&Germany&288\\
\hline\end{tabular}\end{table}
\begin{table}[h]\centering\small\caption{Top organisms, recent isolates.}\begin{tabular}{lrlr}\hline
\textbf{NDM organism}&\textbf{n}&\textbf{OXA-48-like organism}&\textbf{n}\\\hline
Klebsiella pneumoniae&18,066&Klebsiella pneumoniae&7,352\\
E.coli and Shigella&11,688&E.coli and Shigella&4,367\\
Enterobacter hormaechei&3,319&Citrobacter freundii&501\\
Acinetobacter baumannii&2,167&Klebsiella oxytoca&239\\
Providencia alcalifaciens&1,365&Enterobacter hormaechei&122\\
Pseudomonas aeruginosa&980&Acinetobacter baumannii&65\\
Citrobacter freundii&764&Providencia alcalifaciens&54\\
Klebsiella oxytoca&569&Serratia marcescens&24\\
\hline\end{tabular}\end{table}
\section{OXA-48-like family definition}
The 67 family members with amino-acid identity to OXA-48 (global alignment, full length).
\begin{longtable}{lrlr}
\caption{OXA-48-like family membership.}\\
\textbf{allele}&\textbf{identity}&\textbf{allele}&\textbf{identity}\\\hline
\endhead
OXA-48&1.000&OXA-162&0.996\\
OXA-244&0.996&OXA-245&0.996\\
OXA-252&0.996&OXA-370&0.996\\
OXA-505&0.996&OXA-519&0.996\\
OXA-566&0.996&OXA-894&0.996\\
OXA-918&0.996&OXA-920&0.996\\
OXA-933&0.996&OXA-934&0.996\\
OXA-1039&0.996&OXA-1167&0.996\\
OXA-1212&0.996&OXA-1242&0.996\\
OXA-1306&0.996&OXA-204&0.993\\
OXA-514&0.993&OXA-515&0.993\\
OXA-546&0.993&OXA-1038&0.993\\
OXA-1213&0.993&OXA-1304&0.993\\
OXA-793&0.990&OXA-1146&0.990\\
OXA-199&0.989&OXA-416&0.989\\
OXA-538&0.989&OXA-547&0.989\\
OXA-517&0.986&OXA-1055&0.986\\
OXA-181&0.985&OXA-567&0.983\\
OXA-1012&0.982&OXA-405&0.982\\
OXA-1200&0.982&OXA-232&0.981\\
OXA-484&0.981&OXA-923&0.981\\
OXA-924&0.981&OXA-1119&0.981\\
OXA-1181&0.981&OXA-1211&0.981\\
OXA-1226&0.981&OXA-1305&0.981\\
OXA-1308&0.981&OXA-163&0.978\\
OXA-929&0.978&OXA-1240&0.978\\
OXA-833&0.977&OXA-922&0.977\\
OXA-1201&0.977&OXA-1205&0.977\\
OXA-1207&0.977&OXA-1307&0.977\\
OXA-1309&0.977&OXA-438&0.976\\
OXA-247&0.974&OXA-439&0.974\\
OXA-788&0.974&OXA-54&0.924\\
OXA-436&0.913&OXA-535&0.913\\
OXA-731&0.913&&\\
\hline\end{longtable}

\section{Co-carriage combinations}
5,803 assemblies carry both blaNDM and blaOXA-48-like genes. Top combinations:
\begin{longtable}{llr}
\caption{NDM x OXA-48-like co-carriage (assemblies).}\\
\textbf{NDM allele}&\textbf{OXA-48-like allele}&\textbf{assemblies}\\\hline
\endhead
blaNDM-5&blaOXA-181&1,915\\
blaNDM-1&blaOXA-48&1,747\\
blaNDM-1&blaOXA-232&1,482\\
blaNDM-5&blaOXA-48&1,456\\
blaNDM-5&blaOXA-232&1,143\\
blaNDM-1&blaOXA-181&385\\
blaNDM-4&blaOXA-181&332\\
blaNDM-5&blaOXA-1207&119\\
blaNDM-5&blaOXA-244&108\\
blaNDM&blaOXA-48&98\\
blaNDM-5&blaOXA-1205&95\\
blaNDM-5&blaOXA-484&80\\
blaNDM&blaOXA-181&72\\
blaNDM&blaOXA-232&66\\
blaNDM-4&blaOXA-48&48\\
\hline\end{longtable}
\section{Data provenance and query record}
All population data were pulled 2026-09-23 from the NCBI Pathogen Detection isolate backend (the solr API behind the public Isolates Browser), collection ``amr'' (AMRFinderPlus calls; refgene version recorded per row) and collection ``isolates'' (SNP-cluster join, 100-accession batches). Reference data: CARD latest data tarball (4,366,637 bytes, sha256 94b89501fc26daafbe65e4151c17a8b007dc1272c4f8201737e44e44cdd4c065); NCBI beta-lactamase allele registry Allele.tab/Allele-dna.fa (1,557 alleles). Structures: RCSB PDB 3SPU, 4EYL, 4S2P, 6P97 downloaded 2026-09-23. Assemblies: NCBI Datasets v2 API, GENOME\_FASTA, 711 of 722 targets delivered. Software: blastn 2.17.0+, Python 3.10.12, Biopython 1.88. Manifest: ndm-oxa/data/MANIFEST.sha256 (44 payloads).

\section{Gate evidence}
G1: 136/136 reference self-recovery (blastn, CARD reference DB). G2: 1,599/1,607 (99.50\%) vs NCBI AMRFinderPlus; 8 disagreements resolved. G4: spike-in controls detected; published-genome control pending.
\section{Verified/thin/missing summary}
Verified: allele calls (G1/G2-backed), stop-codon porin disruptions, structural distances from PDB coordinates, H250N and L10-loop novel variants (sequence-verified). Thin: porin indel classes (loop-polymorphism confound), OXA-48-like borderline family members, GCA\_059932675.2 L10-loop cluster (possible sequencing artifact). Missing: 11 undelivered assemblies, MIC phenotypes (not in Pathogen Detection), published porin-loss genome control.
\end{document}
